\chapter{The Ocean: Prognostic-Momentum Core}\label{chap:ch07_ocean_prognostic_core}

% --- local symbols (minimal; not in the shared preamble) ---
\providecommand{\taux}{\tau^{x}}                 % zonal wind stress
\providecommand{\tauy}{\tau^{y}}                 % meridional wind stress
\providecommand{\wbc}{\delta_{\mathrm{S}}}        % Stommel western boundary-layer width

The ocean velocities used by the live, coupled \model{} are \emph{diagnostic}: the
horizontal flow is read off the density field through an algebraic thermal-wind inversion,
and the overturning is closed with an interim Stommel/box parameterization
(Chapter~\ref{chap:ch08_ocean_overturning}). That design is numerically robust but it
breaks a property we want for tipping-point science: in the diagnostic core the sensitivity
of the Atlantic overturning to a density perturbation, $\partial(\mathrm{AMOC})/\partial\dens$,
is controlled by the closure rather than by momentum dynamics, so a bifurcation cannot be
\emph{earned} by the physics. This chapter documents the prognostic-momentum ocean core that
is being built to replace that closure. It is a set of standalone, fully differentiable
building blocks --- a variable-coefficient elliptic solver, a prognostic barotropic
vorticity equation, and a prognostic baroclinic momentum equation --- each validated against
an analytic balance. \textbf{None of it is yet wired into \code{veros\_driver}/\code{main};}
it lives apart precisely so that it can be validated with zero regression risk to the running
model. The standard reference for every balance below is \citet{vallis2017}.

\section{The variable-coefficient elliptic invert}

A streamfunction $\psib$ is the natural prognostic object for a rotating, nearly-nondivergent
ocean: the horizontal transport is $\uvec=\kvec\times\grad\psib$, which is divergence-free by
construction. The vorticity of this transport defines $\psib$ through an elliptic equation. For
a flat-bottom ocean this is Poisson's equation $\lap\psib=\vort$; over real bathymetry it
becomes a \emph{variable-coefficient} elliptic problem,
\begin{equation}\label{eq:varcoef-cart}
  \diver\!\big(\,c\,\grad\psib\,\big)=\vort,
  \qquad c=\frac{1}{\dep}.
\end{equation}
The coefficient $c=1/\dep$ encodes topography: where the column is shallow, a given transport
streamfunction implies a larger vorticity, and the steering of flow along $\Cor/\dep$ contours
is exactly the joint-effect-of-baroclinicity-and-relief (JEBAR) mechanism that makes deep
western boundary currents follow the continental slope \citep{vallis2017}. Setting $c=1$
recovers the flat-bottom Poisson limit. The same operator with $c=\mathrm{const}$ also serves
as a Laplacian for a viscous $\nu\lap\vort$ term.

\subsection*{Discrete operator (Cartesian)}

The code discretizes Eq.~\eqref{eq:varcoef-cart} in flux form on a regular grid that is
periodic in $x$ and Dirichlet ($\psib=0$) in $y$ and on land. The coefficient is averaged to
cell faces (east/west/north/south) as the arithmetic mean of the two adjacent cells,
\begin{equation}
  c_{\mathrm E}=\tfrac12\big(c_{i,j}+c_{i,j+1}\big),\quad
  c_{\mathrm W}=\tfrac12\big(c_{i,j}+c_{i,j-1}\big),\quad\text{(and likewise }c_{\mathrm N},c_{\mathrm S}\text{)},
\end{equation}
and the operator is the net face-flux divergence,
\begin{equation}\label{eq:elliptic-op}
  (L\psib)_{i,j}=
  \frac{c_{\mathrm E}(\psib_{i,j+1}-\psib_{i,j})-c_{\mathrm W}(\psib_{i,j}-\psib_{i,j-1})}{\Delta x^{2}}
  +\frac{c_{\mathrm N}(\psib_{i+1,j}-\psib_{i,j})-c_{\mathrm S}(\psib_{i,j}-\psib_{i-1,j})}{\Delta y^{2}}.
\end{equation}
\codref{chronos\_esm/ocean/solver.py}{\code{\_varcoef\_faces}, \code{\_elliptic\_op}, \code{solve\_elliptic\_varcoef}}

Two design choices deserve emphasis because they are easy to get wrong. First, the faces are
\emph{not} masked: an ocean cell adjacent to land still couples to the land cell, which is
pinned to $\psib=0$ by an identity row in the matrix. This realizes the correct Dirichlet
(no-streamfunction) coastal condition. Masking the faces instead would impose a Neumann
(no-flux) condition, which leaves the operator singular up to an additive constant and makes
the iterative solver diverge --- the source comments call this out explicitly. Second, $c=1/\dep$
must be regularized to a finite value on land before the solve, since $\dep\to0$ there.

The system is made symmetric positive-definite by solving $-L\psib=-\vort$ on ocean cells and
the identity $\psib=0$ on land, then inverted with a Jacobi-preconditioned conjugate-gradient
iteration whose diagonal is $(c_{\mathrm E}+c_{\mathrm W})/\Delta x^2+(c_{\mathrm N}+c_{\mathrm S})/\Delta y^2$.
The CG loop is written with \code{jax.lax.scan} (not \code{while\_loop}) so that it is
differentiable in reverse mode --- the solver is exact in $\psib$ and exposes a clean gradient
$\partial\psib/\partial c$, $\partial\psib/\partial\vort$
(Chapter~\ref{chap:ch03_differentiable}).
\codref{chronos\_esm/ocean/solver.py}{\code{solve\_cg}, \code{jacobi\_preconditioner}}

\subsection*{Spherical face-conductance form}

On the lat--lon grid the elliptic operator carries the spherical metric,
\begin{equation}\label{eq:varcoef-sphere}
  \frac{1}{\Erad^{2}\cos\lat}
  \left[\pd{}{\lon}\!\left(\frac{c}{\cos\lat}\pd{\psib}{\lon}\right)
       +\pd{}{\lat}\!\left(c\cos\lat\,\pd{\psib}{\lat}\right)\right]=\vort .
\end{equation}
Discretizing the area-divided form of Eq.~\eqref{eq:varcoef-sphere} directly would yield a
\emph{non-symmetric} matrix, because the cell area $\Erad^{2}\cos\lat\,\Delta\lon\,\Delta\lat$
varies with latitude. The code therefore multiplies through by the cell area and works in
symmetric \emph{face-conductance} form. A conductance is (face coefficient)$\times$(face
length)$/$(centre-to-centre distance); the Earth radius cancels, leaving
\begin{equation}
  C_{\mathrm E}=\bar c_{\mathrm E}\,\frac{\Delta\lat}{\cos\lat\,\Delta\lon},\qquad
  C_{\mathrm N}=\bar c_{\mathrm N}\,\frac{\cos\lat_{\mathrm N}\,\Delta\lon}{\Delta\lat},
\end{equation}
with $\bar c$ the face-averaged coefficient and $\lat_{\mathrm N}$ the north-face latitude
(symmetrically for W, S). The symmetric operator and right-hand side are
\begin{equation}\label{eq:sphere-op}
  (M\psib)_{i,j}=\sum_{\mathrm{faces}}C_{\mathrm{face}}\,(\psib_{i,j}-\psib_{\mathrm{nbr}}),
  \qquad b=-\vort\,\big(\Erad^{2}\cos\lat\,\Delta\lon\,\Delta\lat\big),
\end{equation}
which is exactly $-\,\mathrm{area}\times\diver(c\grad\psib)$ and is SPD, so the same CG applies.
\codref{chronos\_esm/ocean/solver.py}{\code{\_sphere\_conductances}, \code{\_sphere\_op}, \code{solve\_elliptic\_varcoef\_sphere}}

\begin{remark}
The factor $1/\cos\lat$ in the zonal conductance blows up at the pole, the same lat--lon
singularity that destabilized the single-level atmosphere (Chapter~\ref{chap:ch10_atmosphere_singlelevel}).
The fix is identical in spirit: $\cos\lat$ is floored at \code{cos\_min}$=0.087$
($\approx\lat=85^\circ$) so the conductance stays finite. This is a deliberate, quantified
approximation in the polar caps, not a bug.
\end{remark}

A manufactured-solution test confirms the round trip $\psib\!\to\!\vort\!\to\!\psib$ to
relative error $<2\times10^{-3}$ (Cartesian, variable $c$) and $<10^{-4}$ (spherical), and that
the constant-$c$ solver reproduces the plain Poisson solve.
\codref{tests/test\_barotropic\_gyre.py}{manufactured-solution and Poisson-consistency tests}

\section{The prognostic barotropic vorticity equation}

With the invert in place, the depth-integrated (barotropic) circulation is advanced by a
\emph{prognostic} relative-vorticity equation rather than diagnosed. In the linear,
quasi-geostrophic form the model integrates,
\begin{equation}\label{eq:barotropic}
  \pd{\vort}{t}=-\beq\,\pd{\psib}{x}-r\,\vort+F+\nu\lap\vort,
  \qquad \uu=-\pd{\psib}{y},\ \ \vv=\pd{\psib}{x},
\end{equation}
where $\beq=\dd \Cor/\dd y$ is the planetary-vorticity gradient, $r$ a linear bottom drag, $\nu$ a
lateral viscosity, and the forcing is the wind-stress curl
\begin{equation}
  F=\frac{1}{\densz\,\dep}\,\big(\curl\boldsymbol{\tau}\big)_z
   =\frac{1}{\densz\,\dep}\!\left(\pd{\tauy}{x}-\pd{\taux}{y}\right).
\end{equation}
The term $-\beq\,\partial\psib/\partial x=-\beq\vv$ is the advection of planetary vorticity by
the meridional transport --- the source of the $\beta$-effect and of western intensification
\citep{vallis2017}. Each step is forward Euler in $\vort$, after which $\psib$ is re-diagnosed
from $\vort=\diver(c\grad\psib)$ using the elliptic invert of the previous section, warm-started
from the previous $\psib$:
\begin{equation}
  \vort^{\,n+1}=\vort^{\,n}+\Delta t\big(-\beq\,\partial_x\psib^{\,n}-r\vort^{\,n}+F+\nu\lap\vort^{\,n}\big),
  \qquad \psib^{\,n+1}=L^{-1}_{c}\,\vort^{\,n+1}.
\end{equation}
\codref{chronos\_esm/ocean/barotropic.py}{\code{step\_barotropic}, \code{velocities\_from\_psi}, \code{wind\_stress\_curl}}

On the sphere the planetary-vorticity term simplifies elegantly. Since $\beq=2\Om\cos\lat/\Erad$
and $\vv=(\Erad\cos\lat)^{-1}\partial\psib/\partial\lon$, the $\cos\lat$ factors cancel:
\begin{equation}\label{eq:beta-sphere}
  \beq\,\vv=\frac{2\Om\cos\lat}{\Erad}\cdot\frac{1}{\Erad\cos\lat}\pd{\psib}{\lon}
          =\frac{2\Om}{\Erad^{2}}\pd{\psib}{\lon},
\end{equation}
so the spherical step advances $\partial_t\vort=-(2\Om/\Erad^{2})\,\partial_\lon\psib-r\vort+F$
with the spherical curl and metric velocities $\uu=-\Erad^{-1}\partial_\lat\psib$,
$\vv=(\Erad\cos\lat)^{-1}\partial_\lon\psib$.
\codref{chronos\_esm/ocean/barotropic.py}{\code{step\_barotropic\_sphere}, \code{velocities\_sphere}, \code{wind\_stress\_curl\_sphere}}

\subsection*{Analytic validation: the Stommel gyre}

At steady state ($\partial_t\vort=0$, $\nu=0$, $c=1$) Eq.~\eqref{eq:barotropic} reduces to the
Stommel balance \citep{stommel1961},
\begin{equation}\label{eq:stommel}
  r\lap\psib+\beq\,\pd{\psib}{x}=F .
\end{equation}
Away from the western wall the friction term is negligible and the interior obeys the Sverdrup
relation $\beq\,\partial_x\psib=F$, i.e. $\vv=F/\beq$ \citep{sverdrup1947}. The return flow is collected into a thin
western boundary layer of width
\begin{equation}\label{eq:wbc}
  \wbc=\frac{r}{\beq},
\end{equation}
which is what produces western intensification (a strong, narrow Gulf-Stream-like jet on the
west, weak broad return flow in the interior). The validation spins the gyre up from rest under
a single-gyre wind-stress curl and checks both signatures: the meridional jet $|\vv|=|\partial_x\psib|$
peaks in the western third of the basin, and the interior $\partial_x\psib$ matches $F/\beq$ to
better than $35\%$ (Cartesian) and $45\%$ (spherical). With $r=5\times10^{-6}\,\mathrm{s^{-1}}$
and $\beq\approx2\times10^{-11}\,\mathrm{m^{-1}s^{-1}}$ this gives $\wbc\approx250\,\mathrm{km}$, a
few grid cells --- just resolved.
\codref{tests/test\_barotropic\_gyre.py}{\code{test\_stommel\_gyre\_sverdrup\_and\_western\_intensification}, spherical variant}

\begin{example}[Sverdrup transport from the wind]
For the subtropical North Atlantic, a curl of order
$(\curl\boldsymbol\tau)_z\sim-2\times10^{-7}\,\mathrm{N\,m^{-3}}$ over
$\beq\approx2\times10^{-11}\,\mathrm{m^{-1}s^{-1}}$ gives an interior southward transport per unit
width $\vv\dep=(\curl\boldsymbol\tau)_z/(\densz\beq)\sim-10\,\mathrm{m^2\,s^{-1}}$. Integrated across
an ocean width of $\sim6000\,\mathrm{km}$ this is $\mathcal O(30)\Sv$, the right order for the
wind-driven gyre --- and the same $30\Sv$ must return in the western boundary current. The model
reproduces this partition; verify it yourself in Exercise~\ref{ex:sverdrup}.
\end{example}

\section{The prognostic baroclinic momentum equation}

The barotropic core fixes the gyre circulation but not the overturning, which is set by
\emph{density gradients} acting through the vertical structure of the flow. This is the role of
the baroclinic momentum core. The continuous equation, in the hydrostatic, linear-drag form the
code integrates, is
\begin{equation}\label{eq:momentum}
  \pd{\uvec}{t}+\Cor\,\kvec\times\uvec
  =-\frac{1}{\densz}\grad\pres-r\,\uvec+\nu\lap\uvec,
  \qquad \pd{\pres}{z}=-\dens\,\grav ,
\end{equation}
with $\pres$ the hydrostatic pressure. The Coriolis term $\Cor\,\kvec\times\uvec$ is the
load-bearing balance: away from boundaries and friction, Eq.~\eqref{eq:momentum} at steady state
is geostrophy, $\Cor\,\kvec\times\uvec=-\densz^{-1}\grad\pres$, whose vertical derivative is
\emph{thermal wind},
\begin{equation}\label{eq:thermal-wind}
  \Cor\,\pd{\vv}{z}=\frac{\grav}{\densz}\,\pd{\dens}{x},\qquad
  \Cor\,\pd{\uu}{z}=-\frac{\grav}{\densz}\,\pd{\dens}{y}.
\end{equation}
Equation~\eqref{eq:thermal-wind} is precisely how density drives the flow: a horizontal density
gradient is balanced by a vertical shear of the horizontal velocity. The interim diagnostic core
imposes Eq.~\eqref{eq:thermal-wind} algebraically; the prognostic core lets it \emph{emerge} as
the steady limit of Eq.~\eqref{eq:momentum}, so the same density field also drives the
ageostrophic, overturning component through the friction and time-tendency terms.

\subsection*{Discrete pressure-gradient force}

The hydrostatic pressure is the running vertical integral of density from the surface down,
\begin{equation}
  \pres(z)=\grav\!\int_{z}^{0}\!\dens\,\dd z'\ \longrightarrow\
  \pres_{k}=\grav\sum_{k'\le k}\dens_{k'}\,\Delta z_{k'}
  \quad(\text{a }\code{cumsum}\text{ from the surface}),
\end{equation}
defined up to an irrelevant surface constant. The horizontal pressure-gradient acceleration is
$-\densz^{-1}\grad\pres$, with centred differences (periodic in $x$, one-sided at the $y$-edges).
\codref{chronos\_esm/ocean/momentum.py}{\code{hydrostatic\_pressure}, \code{pressure\_gradient\_accel}}

\subsection*{Semi-implicit Coriolis}

The one numerical subtlety is the Coriolis term. An \emph{explicit} treatment of $\Cor\,\kvec\times\uvec$
is unconditionally unstable --- it rotates the velocity vector by a finite angle each step and the
amplitude grows without bound for any $\Delta t$. The code therefore treats Coriolis
\emph{implicitly}, which (with all other forcing collected into $F_u,F_v$) is a $2\times2$ rotation
solve:
\begin{equation}\label{eq:coriolis-impl}
  \frac{\uu'-\uu}{\Delta t}=\Cor\,\vv'+F_u,\quad
  \frac{\vv'-\vv}{\Delta t}=-\Cor\,\uu'+F_v
  \ \Longrightarrow\
  \begin{cases}
    \uu'=\dfrac{r_u+\Delta t\,\Cor\,r_v}{1+(\Delta t\,\Cor)^2},\\[2.2ex]
    \vv'=\dfrac{r_v-\Delta t\,\Cor\,r_u}{1+(\Delta t\,\Cor)^2},
  \end{cases}
\end{equation}
where $r_u=\uu+\Delta t\,F_u$, $r_v=\vv+\Delta t\,F_v$. The denominator $1+(\Delta t\,\Cor)^2$ is the
key: it is $\ge1$ for every $\Delta t$, so the update is a contraction and is stable for arbitrarily
large $\Delta t\,\Cor$. The full step assembles $F_u=a_x-r\uu+\nu\lap\uu$ from the pressure-gradient
acceleration $a_x$, linear drag, and optional viscosity, then applies Eq.~\eqref{eq:coriolis-impl}.
\codref{chronos\_esm/ocean/momentum.py}{\code{coriolis\_semi\_implicit}, \code{step\_momentum}}

\begin{lstlisting}[caption={Semi-implicit Coriolis: the $1+(\Delta t f)^2$ denominator makes it stable for any timestep.}]
ru = u + dt * Fu;  rv = v + dt * Fv
d  = 1.0 + (dt * f)**2
u_new = (ru + dt*f*rv) / d
v_new = (rv - dt*f*ru) / d
\end{lstlisting}

\subsection*{Analytic validation: thermal wind}

The momentum core is spun up from rest under a fixed, horizontally-varying density field and the
emergent vertical shear is compared with the thermal-wind prediction of Eq.~\eqref{eq:thermal-wind}.
In the interior the model shear matches analytic thermal wind to relative error $<12\%$, and a
separate test confirms the semi-implicit Coriolis stays bounded ($|\uu|<5\,\mathrm{m\,s^{-1}}$) at a
timestep that would explode under explicit Coriolis.
\codref{tests/test\_momentum.py}{\code{test\_thermal\_wind\_balance}, \code{test\_coriolis\_semi\_implicit\_stable}}

\begin{proposition}[Density drives a clean overturning]
Couple the baroclinic shear of Eq.~\eqref{eq:thermal-wind} to mass continuity. A meridional
density (hence buoyancy) contrast sets, via thermal wind, a sheared meridional flow; continuity
closes it in the vertical, producing an overturning cell whose strength scales monotonically with
the density contrast. Because the velocities here are \emph{prognostic} solutions of
Eq.~\eqref{eq:momentum}, the map (density)$\to$(overturning) is a genuine, differentiable function
of the dynamics, so $\partial(\mathrm{AMOC})/\partial\dens\neq0$ for the right physical reason ---
the precondition for a clean, rather than closure-imposed, AMOC bifurcation
(Chapter~\ref{chap:ch17_amoc_tipping}). Contrast this with \citet{marotzke1991}-style box models,
where the overturning--density relation is \emph{postulated}.
\end{proposition}

\section{Status: validated standalone, not yet wired in}

To be unambiguous: the three components above (\code{solver.py}, \code{barotropic.py},
\code{momentum.py}) are validated against the Stommel/Sverdrup gyre balance, manufactured elliptic
solutions, and thermal-wind balance, and every piece is differentiable end-to-end (the gyre
amplitude has a finite, nonzero gradient with respect to the wind-stress amplitude; the momentum
solution with respect to the density forcing). They are \textbf{not} called by the live ocean driver
\code{veros\_driver} or by \code{main.step\_coupled}. The coupled model still uses the diagnostic
thermal-wind velocities plus the interim box overturning closure of
Chapter~\ref{chap:ch08_ocean_overturning}. Wiring the prognostic core into the live model ---
replacing the algebraic inversion with the time-integrated barotropic$+$baroclinic split, on the
real bathymetric $c=1/\dep$ over the spherical grid --- is the next major ocean milestone. Until
then this chapter documents a tested, trustworthy laboratory, not a production path.

The differentiability is not incidental. Because both inverts and both time-steppers are written in
pure \jax{} with \code{lax.scan} loops, the entire (density$\to$velocity$\to$transport) chain admits
reverse-mode gradients, which is what will let later chapters fit the core to observations and probe
the tipping manifold (Chapters~\ref{chap:ch15_differentiable_applications},
\ref{chap:ch17_amoc_tipping}). For the GM/Redi eddy parameterization that the prognostic core will
ultimately interact with, see Chapter~\ref{chap:ch05_ocean_mixing} and
\citet{gentmcwilliams1990,griffies2004}.

\section*{Exercises}

\begin{exercise}[Western boundary-layer width]
Using \code{spin\_up\_gyre} from \codref{chronos\_esm/ocean/barotropic.py}{barotropic gyre},
spin up the Stommel gyre for three values of the bottom drag, $r\in\{2.5,5,10\}\times10^{-6}\,\mathrm{s^{-1}}$,
holding $\beq$ fixed. Locate the column of peak $|\vv|=|\partial_x\psib|$ and estimate the boundary-layer
width in each case. Verify that it scales as $\wbc=r/\beq$ (Eq.~\ref{eq:wbc}). At what $r$ does
$\wbc$ shrink below one grid cell, and what happens to the solution then?
\end{exercise}

\begin{exercise}[Sverdrup balance in the interior]\label{ex:sverdrup}
For the spun-up gyre, plot $\partial_x\psib$ along a mid-basin latitude and overlay $F/\beq$.
Confirm the Sverdrup balance holds in the interior and fails inside $\wbc$ of the west wall. By how
much does the basin-integrated interior (Sverdrup) transport differ from the western boundary-current
return transport? They should be equal and opposite; explain any residual in terms of the friction and
discretization.
\end{exercise}

\begin{exercise}[Thermal wind from a front]
With \code{spin\_up} from \codref{chronos\_esm/ocean/momentum.py}{baroclinic momentum}, impose a
density field with a single meridional front (e.g. a $\tanh$ in $y$, uniform in $x$) and integrate to
steady state. Diagnose the vertical shear $\partial\uu/\partial z$ and compare it with
$-\,(\grav/\densz\Cor)\,\partial\dens/\partial y$ from Eq.~\eqref{eq:thermal-wind}. Then double the
front strength and confirm the shear --- and the implied overturning --- doubles. Why is this the
mechanism a box model cannot capture from first principles \citep{marotzke1991}?
\end{exercise}

\begin{exercise}[Why explicit Coriolis explodes]
Disable the semi-implicit solve and step $\uu,\vv$ with an \emph{explicit} Coriolis term,
$\uu'=\uu+\Delta t(\Cor\vv+F_u)$, $\vv'=\vv+\Delta t(-\Cor\uu+F_v)$. Show analytically that the
amplification factor of the unforced rotation is $\sqrt{1+(\Delta t\Cor)^2}>1$, so the kinetic energy
grows every step regardless of $\Delta t$. Then show the semi-implicit update of
Eq.~\eqref{eq:coriolis-impl} has amplification exactly $1$ for the unforced case. Confirm numerically
with \code{test\_coriolis\_semi\_implicit\_stable} as a starting point.
\end{exercise}
